\subsection{Convergence Point: TD-MPC2}
\label{section:fc_tdmpc2}

TD-MPC2 \citep{Hansen2024TDMPC2} is the convergence point of three lines that the preceding sections traced separately. Model predictive control with learned dynamics, in the style of PETS and MBPO (\S\ref{section:fc_mbpo}), supplies the planning frame. Recurrent latent world models, in the Dreamer line (\S\ref{section:fc_rssm}), supply the idea that the model should live in a learned latent and need not decode observations. Value-aware learning, articulated by MuZero and the value-equivalence literature (\S\ref{section:fc_value_aware}), supplies the target the model is trained against. TD-MPC2 keeps the planner of the first line, drops the observation decoder of the second, and adopts the value functional as the training signal of the third. The single empirical commitment that organizes the paper is that one fixed hyperparameter configuration should work across many tasks and many model sizes, with no per-task tuning and no rescaling of components as the network grows.

\subsubsection{Components and joint loss}
\label{section:fc_tdmpc2_components}

The agent learns five components jointly. An encoder $h$ maps an observation $s_t$ to a normalized latent $z_t = h(s_t)$. A latent dynamics model $d$ predicts the next latent, $\hat z_{t+1} = d(z_t, a_t)$. A reward model $R$ predicts the immediate reward from the latent and action, $\hat r_t = R(z_t, a_t)$. A terminal action-value function $Q$ supplies a bootstrap beyond the planning horizon, $\hat q_t = Q(z_t, a_t)$. A policy prior $p$ supplies an amortized action proposal $\hat a_t = p(z_t)$ used both to warm-start the planner and to seed the bootstrap targets. The joint loss is the sum of a latent consistency term that pushes the predicted next latent to agree with the encoded next observation, $\| d(z_t, a_t) - \operatorname{sg}(h(s_{t+1})) \|^2$ under a stop-gradient on the target encoding, a cross-entropy loss on the reward prediction, a cross-entropy loss on the temporal-difference targets for $Q$, and a behavioral cloning term that regresses $p$ on the actions selected by the planner. Notably absent is an observation reconstruction loss; the model is never asked to decode pixels or proprioceptive features and is trained only against quantities the value calculation actually consumes. This is the architectural commitment that distinguishes TD-MPC2 from the Dreamer line (\S\ref{section:fc_rssm}), in which the world model carries a reconstruction term throughout.

\subsubsection{MPPI planning with value bootstrap}
\label{section:fc_tdmpc2_mppi}

Action selection at each environment step uses Model Predictive Path Integral control, a sampling-based optimizer over Gaussian action sequences \citep{Hansen2024TDMPC2}. The planner maintains a sequence of Gaussian distributions $(\mathcal{N}(\mu_h, \sigma_h^2))_{h=t}^{t+H-1}$ over the next $H$ actions, draws a population of candidate sequences, rolls each one forward through the learned latent dynamics, scores the imagined returns, and refits $(\mu_h, \sigma_h)$ to a softmax-weighted average of the top-scoring sequences. The scoring rule is the finite-horizon return augmented with a terminal value bootstrap,
\[
  \mu^\star, \sigma^\star \;=\; \arg\max_{(\mu,\sigma)} \;\mathbb{E}_{(a_t, \ldots, a_{t+H}) \sim \mathcal{N}(\mu, \sigma^2)} \Big[\, \gamma^{H} Q(z_{t+H}, a_{t+H}) \;+\; \sum_{h=t}^{H-1} \gamma^{h-t} R(z_h, a_h) \,\Big],
\]
where the trajectory expectation is taken under the learned latent dynamics $d$ rolled from the current encoded state $z_t = h(s_t)$. The terminal $Q$ extends the planner's effective horizon past the finite window $H$, so the agent does not need to roll long open-loop sequences through the learned model to capture long-horizon credit. This is the architectural reason TD-MPC2 sidesteps the long-rollout compounding error that limits PETS (\S\ref{section:fc_mbpo}), which has no terminal value and must extend $H$ to capture distant reward. After the planner converges, the agent executes only $a_t \sim \mathcal{N}(\mu_t^\star, (\sigma_t^\star)^2)$ and re-plans at the next step. The amortized policy $p$ seeds the next planning call and supplies the action used in the $Q$-bootstrap.

\subsubsection{Empirical headline and scaling}
\label{section:fc_tdmpc2_empirics}

The paper evaluates TD-MPC2 across 104 continuous-control tasks spanning four benchmarks, DeepMind Control, Meta-World, ManiSkill2, and MyoSuite, under a single fixed hyperparameter configuration with no per-task tuning. The task suite includes action dimensionalities up to $39$, sparse rewards, multi-object manipulation, musculoskeletal motor control, and locomotion on dog and humanoid embodiments. Across all four benchmarks TD-MPC2 outperforms SAC, the original TD-MPC, and DreamerV3 at comparable parameter counts. The contribution is not that any single task is solved better than ever before, but that one configuration solves all of them, which is the property the field had previously assigned to model-free baselines under heavy per-task tuning.

The scaling experiment trains a single multi-task agent on $80$ tasks drawn from DeepMind Control and Meta-World and varies the world model size from $1$ million parameters to $317$ million. The normalized score moves from $16.0$ at $1$M parameters to $70.6$ at $317$M, roughly log-linear in parameter count over the two and a half orders of magnitude examined, with no observed saturation at the upper end. This is the first model-based reinforcement learning paper to report a clean monotone scaling curve over this range. The same hyperparameters are used at every model size; no learning rate, batch size, or capacity-specific schedule is retuned. The empirical claim of TD-MPC2 is that an MPC-shaped model-based agent, built from the components developed in the three earlier lines, scales with parameters in the same way that supervised and self-supervised models do.
